\documentclass[../../main.tex]{subfiles}

\graphicspath{{../../image/}{../../figure/}}

\begin{document}

% 原书物理页 10--14，印刷页码 XII--XV。
% 原书为无编号的卷首篇章；除 (HCL) 一处术语定义外，通篇为概述性叙述，故保留正文段落，
% 仅将该术语定义按 v13 §3.0 A 整理为无编号 definition*（原书无编号，且此处位于章计数为 0 的卷首，
% 有编号环境会显示为“定义 0.1”，故用星号环境）。

\chapter*{引言}
\phantomsection\label{chapter:front-introduction}
\pdfbookmark[1]{引言}{bm-front-introduction}
\markboth{引言}{引言}

连续介质物理的种子是由十八世纪自然哲学家的著作播下的，其中最著名的是 Euler；到十九世纪中叶，树木已完全长成，准备结果。正是在这样的环境中，气体动力学的研究催生了散度形式的拟线性双曲组理论，即通常所谓的双曲守恒律；而这两个主题在过去两百年间一直携手前行。本书意在从双曲守恒律理论与连续介质物理的亲缘关系这一立场出发，呈现前者。\hyperref[chapter:front-history]{双曲守恒律早期历史概述}紧随本引言之后。尽管研究仍以轻快的步伐推进，这两个领域如今都已达到足以支撑这样一部论述的成熟程度。

在连续介质物理的范围内，物体被实现为连续介质，诸如质量、动量与能量等所谓“广延量”，通过其密度场来监控，而这些密度场由平衡律与本构方程相联系。\cref{chapter:ch02-连续介质物理导引}给出了对经典物理这一分支的一个虽嫌简略却自足的介绍。读者可借助该主题的专门教材加以充实。

就原始表述而言，典型的平衡律断言：物体任一子域内所储存的某一广延量的总量的时间变化率，由该量通过该子域边界的通量率与它在该子域内的生成率之和所平衡。若无生成，则平衡的广延量守恒。使连续介质物理适于分析处理的独特之处在于：在相当自然的假设下，上述那种总体平衡的陈述可归结为场方程，即散度形式的偏微分方程。

所采用的平衡律的清点划定并标识出特定的连续介质理论，例如力学、热力学连续介质力学、电动力学等等。在一种连续介质理论的语境中，本构方程编码介质的材料性质，例如导热黏性流体、弹性固体、弹性电介质等。这些本构关系与场方程的耦合，产生了封闭的偏微分方程组，被称为“平衡律”或“守恒律”，连续介质的平衡态或运动即由之确定。从历史上看，绝大多数值得注意的偏微分方程都是通过这一过程产生的。对双曲守恒律方程组而言尤其如此，这从历史叙述中可以见出。本书的中心论题是：连接连续介质物理与偏微分方程理论的脐带不应被剪断，因为它仍在双向输送养分。

平衡律方程组可以是椭圆的，通常在静力学中；可以是双曲的，在动力学中、对具有“弹性”响应的介质；可以是椭圆--双曲混合的，在介质发生相变的静力学或动力学情形；也可以是抛物或抛物--双曲混合的，在存在黏性、热传导或其他扩散机制时。据此，基本概念将在\cref{chapter:ch01-平衡律}中以足以涵盖上述一切可能性的普遍程度引入。尽管如此，由于本书的主题是双曲守恒律，讨论最终将聚焦于这类方程组，自\cref{chapter:ch03-双曲平衡律方程组}开始。

\begin{definition*}[齐次双曲守恒律]
所谓\emph{齐次双曲守恒律}，指的是如下形式的、属于双曲型的一阶偏微分方程组：
\begin{align*}
(\mathrm{HCL})\qquad \partial_t H(U)+\sum_{\alpha=1}^{m}\partial_\alpha G_\alpha(U)=0.
\end{align*}
状态向量 $U$ 取值于 $\mathbb{R}^{n}$，待确定为空间变量 $(x_1,\ldots,x_m)$ 与时间 $t$ 的函数。给定函数 $H$ 与 $G_1,\ldots,G_m$ 是 $\mathbb{R}^{n}$ 到 $\mathbb{R}^{n}$ 的光滑映射。符号 $\partial_t$ 表示 $\partial/\partial t$，$\partial_\alpha$ 表示 $\partial/\partial x_\alpha$。双曲性的概念将在\cref{section:ch03-双曲性}中确切说明。
\end{definition*}

双曲守恒律的解可以设想为传播的波。当方程组非线性时，压缩波的剖面会逐渐变陡，最终破裂，产生跳跃间断，这些间断作为激波继续传播。因此，理论不可避免地必须处理弱解。这一困难因如下事实而进一步加剧：在弱解的语境中，唯一性丧失。于是，有必要设计恰当的标准以挑选出可容许的弱解。连续介质物理通过热力学第二定律自然地诱导出这类可容许性标准。这些标准可以被纳入分析理论：或者直接地，径直规定可容许解应满足“熵”不等式；或者间接地，给方程组配上极小量的扩散，这对光滑解的影响可以忽略，但在激波出现时会作出强烈反应，把不满足热力学可容许性的解剔除。在本书中扮演核心角色的“熵”与“消失扩散”这两个概念，将在\cref{chapter:ch03-双曲平衡律方程组}与\cref{chapter:ch04-cauchy问题}中引入。

\cref{chapter:ch05-熵与经典解稳定性}在经典解的语境中讨论双曲平衡律方程组的 Cauchy 问题与初边值问题。结果表明，这些问题局部适定，所得的光滑解是稳定的——甚至在更广泛的可容许弱解类中也是稳定的——但其存在时间有限，除非存在阻止波破裂的耗散源。这一分析凸显了热力学第二定律的稳定化效应，以及刻画松弛的耗散所起的作用。

大范围的 Cauchy 问题只能在弱解的语境中考虑。对多空间维中多于一个方程的方程组而言，这仍是一片未知领域（terra incognita），因为分析目前面临着看似不可逾越的障碍。Cauchy 问题可能根本不是普遍适定的：或者由于唯一性发生灾难性失效（见\cref{section:ch04-Euler方程}），或者由于分布解不存在。在后一情形，人们将不得不诉诸更弱的、取值于测度的解类（见\cref{section:ch17-测度值解与补偿紧致}）。甚至可以设想，双曲组应当被视为具有极小黏性或色散的扩散系统在 Plato 意义上的影子。尽管如此，本书仍将聚焦于成功的例子，即容许标准分布弱解的问题。这些包括一维或多维空间中的标量守恒律、单一空间维中的双曲守恒律方程组，以及当不变性（球对称、定常性、自相似性等）把自变量数目归结为两个时的多空间维方程组。

\cref{chapter:ch06-标量守恒律L1理论}详细呈现了多空间维标量守恒律的 Cauchy 问题与初边值问题的 $L^{\infty}$ 与 $BV$ 解这一丰富而完备的理论。

自\cref{chapter:ch07-一维双曲平衡律组}起，考察的焦点固定于一维空间中的守恒律方程组。在这一框架下，理论具有若干对分析者极有帮助的特殊性质，因而取得了重大进展。

\cref{chapter:ch08-容许激波}系统阐述激波的性质。特别地，引入多种激波可容许性准则，并对它们进行比较与对照。随后在\cref{chapter:ch09-容许波扇与riemann问题}中，把可容许激波与另一类称为中心稀疏波的特殊解相结合，以合成求解经典 Riemann 问题的波扇。Riemann 问题的解又可作为构件，用于在 $BV$ 类（有界变差函数类）中构造 Cauchy 问题的解。本书将给出基于这一途径的两种构造方法：\cref{chapter:ch13-随机选择法}中的随机选择格式，与\cref{chapter:ch14-前沿追踪与标准riemann半群}中的前沿追踪算法。这些解的唯一性与稳定性也将被建立。

\cref{chapter:ch15-消失黏性构造bv解}概述了用消失黏性方法，对一般严格双曲守恒律方程组的 Cauchy 问题构造 $BV$ 解的另一种途径。

\cref{chapter:ch16-平衡律组bv解}讨论用随机选择法，为带耗散源、支配松弛现象的严格双曲平衡律方程组构造 $BV$ 解。

上述构造方法一般适用于初值总变差充分小的情形。这一限制似乎是普遍必要的，因为在某些方程组中，当初值“很大”时，连 Cauchy 问题的弱解都可能在有限时间内爆破。这类灾难是否会发生在连续介质物理场方程的解上，目前是一个重大的未决问题。对一类有限的方程组——不过其中包含若干重要的代表——可用补偿紧致的泛函分析方法构造大初值解。这一途径立足于测度值解与 Young 测度的概念，将在\cref{chapter:ch17-补偿紧致性}中概述。

除存在性与唯一性之外，弱解还有其他值得注意的性质。在\cref{chapter:ch10-广义特征线}中，特征线的概念从经典解推广到弱解；它被用来精确刻画标量守恒律解的正则性与长时间行为（见\cref{chapter:ch11-一维标量守恒律}），以及两个守恒律的方程组解的正则性与长时间行为（见\cref{chapter:ch12-两方程真非线性组}）。

最后的\cref{chapter:ch18-多维定常与自相似解}讨论二维空间加时间中几何与不变性把变量数目归结为两个的问题，即标量守恒律的 Riemann 问题、气体动力学中的绕流与激波碰撞、弹性动力学中的空化，以及微分几何中曲面的等距浸入。

参考文献收录近两千条，相当广博却远非全面。除近期进展外，它也提供了对可能已被取代的早期工作的指引，以便读者追踪该领域的演进。

为凸显基本思想，讨论由一般走向特殊，尽管相反的次序显然有其教学上的优点。即便如此，在适当的指导下，本书也可用作教科书。有鉴于此，证明的节奏被有意安排得并不均匀：对可能构成入门课程材料的基本、初等命题放慢；对面向有经验分析者的较高级技术性结果则加快。尽管本书各部分拼合为一个整体，读者仍可在书中选取若干条独立的路线。因此，主要关心双曲守恒律理论与连续介质物理之关联的概念基础者，只需通读\crefrange{chapter:ch01-平衡律}{chapter:ch05-熵与经典解稳定性}。关于标量守恒律的\cref{chapter:ch06-标量守恒律L1理论}几乎可以独立于其余部分阅读。打算把解作为相互作用初等波的复合物来研究的学生，可从\crefrange{chapter:ch07-一维双曲平衡律组}{chapter:ch09-容许波扇与riemann问题}开始，然后或继续读\crefrange{chapter:ch10-广义特征线}{chapter:ch12-两方程真非线性组}，或直接转到\cref{chapter:ch13-随机选择法}、\cref{chapter:ch14-前沿追踪与标准riemann半群}与\cref{chapter:ch16-平衡律组bv解}。类似地，\cref{chapter:ch15-消失黏性构造bv解}只依赖于\cref{chapter:ch07-一维双曲平衡律组}与\cref{chapter:ch08-容许激波}，而\cref{chapter:ch18-多维定常与自相似解}依赖于\cref{chapter:ch03-双曲平衡律方程组}、\cref{chapter:ch07-一维双曲平衡律组}、\cref{chapter:ch08-容许激波}与\cref{chapter:ch09-容许波扇与riemann问题}。最后，\cref{chapter:ch17-补偿紧致性}中所阐述的泛函分析途径只以\cref{chapter:ch07-一维双曲平衡律组}为先修条件。

某些论题或许讨论得过于详细，因为作者对它们特别感兴趣；还有一些结果是在本书中首次发表。另一方面，理论的若干重要方面及其应用则仅被略加触及，或只是极简地勾勒。它们包括多空间维激波与边界条件的稳定性理论、新近兴起的随机初值双曲守恒律理论、由气体动理论导出连续介质物理的平衡律、相变研究以及大量各色应用。其中每一个领域都值得写成一部专门论著。最明显的缺项是对数值方法的讨论——数值方法除其实用价值外，也为理论提供了宝贵的洞见。所幸目前已有若干关于双曲守恒律数值分析的教材可供填补这一空白。

几何测度论、泛函分析与动力系统为双曲守恒律理论提供了必要的工具，但在很大程度上，分析所采用的是量身定制的技巧，带有强烈的几何色彩，强调波的传播与波的相互作用。这或许会给人这样的印象：该领域是一座孤岛，与偏微分方程的大陆相隔绝。然而读者很快就会意识到，双曲守恒律领域影响深远且高度多样，因为它通过桥梁与椭圆方程、抛物方程、色散方程以及动理论方程的疆域相连。

\end{document}